\section{Critério de escolha e recomendação}
\label{sec:criterio}

Aplica-se o critério da \PUm{Secção~7.5} com os parâmetros da
Secção~\ref{sec:parametros}. \novo{A proteção de uma configuração $c$ é a fração
ponderada dos ataques que impede, e o custo agrega as métricas normalizadas por
K0 (\PUm{Eq.~10 e~11}):}\marcaNovo{equações novas}
\begin{equation}
  \begin{gathered}
  S(c) = \frac{\sum_{a \in \mathcal{A}} \omega_a\, \pi_a(c)}{\sum_{a \in \mathcal{A}} \omega_a},
  \\[4pt]
  \operatorname{custo}(c) = \sum_{m} v_m\, r_m(c), \qquad
  r_m(c) = \begin{cases}
    m(c) / m(\mathrm{K0}) & \text{se maior é pior,}\\
    m(\mathrm{K0}) / m(c) & \text{no débito,}
  \end{cases}
  \end{gathered}
  \label{eq:protecao-custo}
\end{equation}
\novo{com os pesos $\omega_a$ e $v_m$ da Secção~\ref{sec:parametros}. Entre as
configurações exequíveis que cumprem os requisitos do caso de uso, as não
dominadas formam a fronteira de Pareto $\mathcal{P}$, e a configuração recomendada
é a mais barata das que atingem a proteção máxima (\PUm{Eq.~12}):}
\begin{equation}
  c^{*} = \operatorname*{arg\,min}_{c \in \mathcal{P},\ S(c) = S_{\max}} \operatorname{custo}(c),
  \qquad S_{\max} = \max_{c \in \mathcal{P}} S(c).
  \label{eq:melhor}
\end{equation}
\novo{Ao longo da fronteira, reporta-se o custo marginal $\Delta\operatorname{custo} / \Delta S$.}

\novo{A Tabela~\ref{tab:criterio} dá, para cada configuração, $S(c)$, o custo, a
exequibilidade e a pertença à fronteira de Pareto em cada caso de uso, e a
Figura~\ref{fig:pareto} mostra a fronteira da telemetria sensível.}

\gerado{pi5/criterion}

\figuraGerada{pi5/fig-pareto}{Proteção $S(c)$ contra custo$(c)$; a tracejado, a
  fronteira de Pareto para a telemetria sensível. \novo{K9, com custo 5,91, está
  representada no limite direito do eixo.}}{fig:pareto}

Os requisitos de cada caso de uso (\PUm{Secção~7.5}) excluem logo K0 a K4: a
telemetria exige integridade, ligação ao tópico e anti-repetição no conteúdo, e
as ordens exigem ainda validade, restrição chave--tópico e prova de origem.

\begin{itemize}[leftmargin=*]
  \item \textbf{Telemetria não sensível e telemetria sensível.} As configurações
    admissíveis são as mesmas nos dois casos (K5 a K8, K10 e K11; a
    confidencialidade exigida pela telemetria sensível é dada pelo modo de cifra,
    que todas usam). K9 cumpre os requisitos da telemetria não sensível mas não é
    exequível: sustenta 5\,000~mensagens/s, metade de $R_{\text{alvo}}$. A
    fronteira de Pareto tem três pontos: K10 ($S = 0{,}64$, custo 1,40), K5
    (0,73; 1,76) e K8 (1,00; 1,80). O custo marginal ao longo da fronteira é de
    3,9 por unidade de $S$ de K10 para K5 e de apenas 0,14 de K5 para K8: os
    últimos 27 pontos de proteção (validade, \novo{restrição chave--tópico,}
    limitação de tráfego e limites do \emph{broker}) quase não custam. $c^* = \text{K8}$ nos dois casos.
  \item \textbf{Ordens a atuadores.} Só K9 dá prova de origem, e é exequível
    para as ordens ($R_{\text{alvo}} = 1\,000$~mensagens/s, que sustenta cinco
    vezes, com p99 de 0,7~ms à carga nominal). $c^* = \text{K9}$, sem
    alternativa: o seu custo, 5,9 vezes o de K0, é o preço da prova de origem.
\end{itemize}

A recomendação antecipada na Parte~1 mantém-se: \textbf{K8 para a telemetria e
K9 para as ordens}. \novo{K11 (ChaCha20) tem a mesma proteção que K8 e custa mais
2\,\% neste processador, que tem AES em \emph{hardware}. Num processador sem essas
extensões, o envelope AES custa 1,4 vezes o do ChaCha20 a 256~octetos, 4 vezes a
4~KiB e 5 a 64~KiB (Secção~\ref{sec:aceleracao}). Aí, K11 passa a ser a escolha
para a telemetria, como a Parte~1 previa.}

\paragraph{Sensibilidade.} A escolha não depende dos pesos fixados. Com todos os
ataques a pesar~1, com pesos de custo iguais, ou com todo o peso no processador
ou na latência p99, $c^*$ continua a ser K8 para a telemetria e K9 para as
ordens. Só com todo o peso nos octetos por mensagem ou no débito máximo é que K8 e
K11 empatam, porque têm o mesmo envelope e saturam no mesmo patamar. A ordem das
configurações com $S = 1$ é sempre a mesma, porque K8 não é pior do que K11 em
nenhuma métrica de custo, e K9 é a mais cara em processador, latência, octetos e
débito.
